\documentclass[11pt, a4paper]{article}

\usepackage[T1]{fontenc}
\usepackage[utf8]{inputenc}
\usepackage{tgpagella}
\usepackage[spanish, es-noshorthands]{babel}
\usepackage{amsmath, amssymb, bm}
\usepackage{booktabs}
\usepackage{tabularx}
\usepackage[table, xcdraw]{xcolor}
\usepackage[most]{tcolorbox}
\usepackage[american]{circuitikz}
\usepackage[margin=2.5cm]{geometry}
\usepackage{fancyhdr}

\pagestyle{fancy}
\fancyhf{}
\setlength{\headheight}{16pt}
\lhead{\textbf{UCB Sede Tarija} | Ing. Mecatrónica}
\chead{}
\rhead{IMT-121 Circuitos Electrónicos I | Diagnóstico S06}
\lfoot{Prof: M.Sc. Ing. Bernardo Quiroga Turdera}
\rfoot{Página \thepage}
\renewcommand{\headrulewidth}{0.4pt}

\definecolor{ucbblue}{RGB}{0, 51, 102}

\begin{document}

\begin{tcolorbox}[colback=ucbblue!5!white, colframe=ucbblue, title=\textbf{Diagnóstico de Punto de Decisión: Superposición}]
    \centering \textbf{Instrucciones:} Resuelve el siguiente problema de manera individual y autónoma (Tiempo objetivo: 10-12 min). Esta actividad no tiene calificación sumativa.
\end{tcolorbox}

\section*{Circuito de Análisis}
\begin{center}
\begin{circuitikz}[scale=1, transform shape, thick]
    \draw (0,0) to[V, l={$V_s = 10\,\text{V}$}, invert] (0,2) to[R, l={$R_1 = 2\,\text{k}\Omega$}] (3,2);
    \draw (3,2) node[above]{$V_o$} to[short, *-] (3,2);
    \draw (3,2) to[R, l_={$R_2 = 3\,\text{k}\Omega$}] (3,0);
    \draw (5,2) to[I, l={$I_s = 1\,\text{mA}$}] (5,0); % De arriba hacia abajo (hacia tierra)
    \draw (3,2) -- (5,2);
    \draw (0,0) -- (5,0);
    \draw (2.5,0) node[ground]{};
\end{circuitikz}
\end{center}

\vspace{0.3cm}
\noindent\textbf{Tarea:} Calcule la tensión $V_o$ aplicando el Teorema de Superposición. \\
Debe incluir de forma explícita los siguientes 6 pasos:

\begin{enumerate}
    \item \textbf{Referencia:} Indique claramente cuál es su terminal positivo y negativo para medir $V_o$.
    \vspace{1.5cm}
    \item \textbf{Subcircuito 1 (Solo $V_s$ activa):} Dibuje el circuito equivalente y calcule la contribución $V_o^{(1)}$.
    \vspace{4.5cm}
    \item \textbf{Subcircuito 2 (Solo $I_s$ activa):} Dibuje el circuito equivalente y calcule la contribución $V_o^{(2)}$. \textit{(Cuidado con el signo)}.
    \vspace{4.5cm}
    \item \textbf{Cálculo de contribuciones con signo:}
    $V_o^{(1)} = $ \rule{3cm}{0.4pt} \hspace{2cm} $V_o^{(2)} = $ \rule{3cm}{0.4pt}
    
    \item \textbf{Suma final:} 
    $V_o = $ \rule{6cm}{0.4pt}
    
    \item \textbf{Justificación conceptual:} Complete la siguiente frase justificando el procedimiento:\\
    \textit{"El teorema de superposición funciona en este circuito porque..."}
    \vspace{1.5cm}
\end{enumerate}

\end{document}
